\documentclass{article}
\usepackage[margin=1in]{geometry}
\usepackage{mathtools, amsfonts, amsthm, xcolor, listings}
\usepackage{hyperref}

\newcommand{\R}{\mathbb{R}}
\newcommand{\N}{\mathbb{N}}
\newcommand{\I}{\mathbb{I}}
\newcommand{\Q}{\mathbb{Q}}
\newcommand{\Z}{\mathbb{Z}}

\newcommand{\set}[1]{\left\{#1\right\}}
\newcommand{\ang}[1]{\left\langle #1 \right\rangle}
\newcommand{\inv}{^{-1}}

\newcommand{\normsub}{\lhd}
\newcommand{\normsube}{\unlhd}

\newcommand{\reflection}[5]{
    \begin{itemize}
    \item Points attempted: #1
    \item Self-assessed score: #2/#1.
    \item Time spent: #3 minutes.
    \item Difficulty: #4/10.
    \item Questions/Feedback: #5.
    \end{itemize}
}


\title{Title}
\author{Your name here}

\begin{document}
\maketitle

\section*{Exercise 2}
\reflection{2}{2}{5}{1}{None.}
Go through the previous homework assignments (including this one) spending at least a minute browsing through each one. Make a note of one or two problems that stood out for you this term, which you would recommend to future students taking this class, and write a note about why. These might be problems where something "clicked" for you, where it helped you to make a connection to new ideas, where you generated a diagram that was meaningful, etc.
  
Some problems that stuck out to me were problems that related back to Linear Algebra.
For example, HW3 Ex3.7 and HW4 Ex2.2.d. These problems stuck out to me because 
it helped me relate abstract algebra and group theory to another area of math 
that I was a little more familiar with at the time. This allowed me to gain a 
strong ``foothold'' on the topic, and helped the topics click for me.
\pagebreak

\section*{Exercise 3}
\reflection{6}{6}{65}{7}{Correctness, Structure}

Note: I ended up scrapping most of my blog, and pivoted topics.
\pagebreak

\section*{Choice Exercise 4}
\reflection{4}{4}{25}{6}{None.}

\href{https://math.stackexchange.com/questions/267521/is-there-a-characterization-of-groups-with-the-property-forall-n-unlhd-g-ex}{
    Is there a characterization of groups with the property \(\forall N \normsube G, \exists H \le G \text{ s.t } H \cong G/N\)?
}

This post really stuck out to me because upon reading the title, I was convinced
that all groups satisfies this property, yet the mere existence of such a post
implies that there exists some groups that don't satisfy this property. I was 
properly baffled! 

As I read through the answers, I wasn't quite able to follow the arguments provided,
but I was definitely intrigued to learn more. Eventually, the post had completely 
lost me, but this is definitely an interesting direction, especially since the 
answer itself admits that there are some ``open questions''.
\pagebreak

\section*{Choice Exercise 8}
\reflection{10}{10}{70}{7/10}{Correctness.}
We call a group action faithful if for all nonidentity elements \( g \in G \setminus \{e\} \), there exists \( x \in X \) such that \( g \cdot x \neq x \). Informally, every non-identity element does something.
\begin{enumerate}
    \item {
        [3] Let \(X = \{1, 2, 3, 4, 5\}\) and let \(G = C_6\). Give an example of a faithful group action of \(G\) onto \(X\). (Hint: is there a one-to-one homomorphism \(\phi\colon C_6 \to S_5\)?)

        Let \(\pi = (123) \in S_5\). We can define a faithful group action of
        \(G = C_6\), by define a homomorphism \(\phi \colon C_6 \to S_5\). Let
        us first enumerate the elements of \(C_6\) as integers from 0 to 5 inclusive.
        We define
        \[
        \phi = \begin{cases}
            \pi^{x/2} &x \text{ is even} \\
            (12) \pi^{(x-1)/2} &x \text{ is odd}
        \end{cases}.
        \]

        This is a ono-to-one homomorphism.
    }

    \item {
        [2] Let \( (G, *) \) be a group that acts on a set \( X \) and let \[         N = \{ g \in G \mid g \cdot x = x \text{ for all } x \in X \}.       \] In class, we showed that \( N \) is a normal subgroup. Repeat the argument here.

        \begin{proof}
            We begin by establishing that \(N\) is a group. First, associativity
            is inherited from \(G\).  Also, by definition of a group action, 
            \(e_G \in N\), since \(e_G \cdot x = x\).
            
            Suppose we have \(n, m \in N\). By definition,
            \(n \cdot x = x\) and \(m \cdot x = x\) for all \(x \in X\). So, 
            by defintion of group actions, \(n \cdot (m \cdot x) = (n * m) \cdot x = x\).
            Because \((n * m) \cdot x = x\), \(n*m \in N\), so \(N\) is closed under \(*\).

            Now, because we know that \(e_G \in N\), we can pick some \(n \in N\), 
            and show that \(n\inv \in N\) because 
            \begin{align}
                e_G \cdot x &= x \\
                (n\inv * n) \cdot x &= \\
                n\inv \cdot (n \cdot x) &= \\
                n\inv \cdot x &= x.
            \end{align}

            Thus, \(N < G\).

            To prove that \(N \normsub G\), we must show that for all \(g \in G\), 
            and \(n \in N\), \(g * n *g\inv \in N\). We show this by
            \begin{align}
                &(g*n*g\inv) \cdot x  \\
                &= (g*n) \cdot g\inv \cdot x \\
                &= g \cdot (n \cdot (g\inv \cdot x)).
            \end{align}

            Now, because \(n \cdot x = x\) for all \(x \in X\), we can simplify to 
            \begin{align}
                &= g \cdot g\inv \cdot x)\\
                &= e_G \cdot x = x.
            \end{align}

            Thus, \(g * n * g\inv \in N\), so \(N \normsub G\).
        \end{proof}
    }

    \item {
        [3] We care about normal subgroups because we can quotient by them. Show that \( G/ N \) acts on \( X \) (denoted \( \bullet \)) by
      \[         gN \bullet x = g \cdot x,       \] and that this map is well-defined: if \(g_1N = g_2N\) then \(g_1 \cdot x = g_2 \cdot x\).
    
        First, notice that \(n \cdot x = x\). Thus, for every \(h \in gN\), 
        \(h = g * n\) for some \(n \in N\). Therefore, \(gN \bullet x = (g * n) \cdot x\) 
        for some \(n \in N\). Thus, \(gN \bullet x = g \cdot x\).

        Naturally, we can then write \(g_1N = g_1 * n_1\) and \(g_2N = g_2 * n_2\) 
        for some \(n_1, n_2 \in N\). So, by substituting, we obtain
        \begin{align}
            g_1N &= g_2N \\
            g_1 * n_1 &= g_2 * n_2. 
        \end{align}

        Now, we apply the group action \(g_1N \bullet x\) and \(g_2N \bullet x\).
        \begin{align}
            g_1N \bullet x &= g_2N \bullet x\\
            (g_1 * n_1) \cdot x &= (g_2*n_2) \cdot x\\
            g_1 \cdot x &= g_2 \cdot x.
        \end{align}
    }

    \item {
        [2] Prove that the action of \( G/ N \) is faithful.

        \begin{proof}
            We begin by restating the faithfulness condition: a group 
            action of \(H\) is faithful if for all \(h \in H\), \(x \in X\), 
            \(h \cdot x = x\) implies \(h = e\).

            From above, we define the group action of \(G/N\) on \(X\) as
            \(gN \bullet x = g \cdot x\).

            By definition, the identity element of \(G/N\) is the coset \(N\).
            Now, suppose we have some \(gN \in G/N\), such that \(gN \bullet x = x\).
            Thus, \(g \cdot x = x\), so \(g \in N\), and \(gN = N\). Therefore, 
            for all \(gN \in G/N\), if \(gN \bullet x = x\), then \(gN\) is the 
            identity element, so this group action must be faithful.
        \end{proof}
    }
\end{enumerate}

\end{document}
